\section{Tiên đề hợp}

\

Như đã khẳng định, nếu $x$ là một tập hợp thì bên trong $x$ có thể chứa các tập hợp khác. Tiên đề hợp phát biểu rằng tồn tại một tập hợp mà nó chứa tất cả và chỉ chứa các phần tử của các tập hợp trong $x$. Viết dưới dạng kí hiệu:
\begin{equation*}
   \defMath{\forall x, \exists u, \forall y \left( y \in u \iff \exists z (z \in x \land y \in z) \right)}.
\end{equation*}

\begin{figure}[H]
   \begin{center}
      \begin{tikzpicture}[>=stealth]
         \draw[thick] (-3,0) ellipse (2.5cm and 3cm);
         \node at (-3, 3.4) {Tập hợp $x = \{A, B\}$};

         % Vẽ tập con X bên trong A
         \draw[thick, dashed] (-3.8, 0.8) ellipse (0.8cm and 1.2cm);
         \node at (-3.8, 2.2) {$A$};

         % Vẽ tập con Y bên trong A
         \draw[thick, dashed] (-2.2, -0.8) ellipse (0.8cm and 1.2cm);
         \node at (-2.2, -2.2) {$B$};

         % Các phần tử trong tập X
         \node[circle, fill, inner sep=1.2pt, label=left:$z_1$] (z1_left) at (-3.8, 1.3) {};
         \node[circle, fill, inner sep=1.2pt, label=left:$z_2$] (z2_left) at (-3.8, 0.3) {};

         % Các phần tử trong tập Y
         \node[circle, fill, inner sep=1.2pt, label=left:$z_3$] (z3_left) at (-2.2, -0.3) {};
         \node[circle, fill, inner sep=1.2pt, label=left:$z_4$] (z4_left) at (-2.2, -1.3) {};

         % --- MŨI TÊN CHUYỂN ĐỔI ---
         \draw[->, ultra thick, colorEmphasis] (0,0) -- (2,0) node[midway, above] {$\bigcup x$};

         % --- VẾ PHẢI: TẬP HỢP \bigcup A ---
         % Vẽ tập hợp kết quả
         \draw[thick] (4.5,0) ellipse (2cm and 2.8cm);
         \node at (4.5, 3.2) {Tập hợp $\bigcup x$};

         % Các phần tử được "giải phóng" ra tập lớn
         \node[circle, fill, inner sep=1.2pt, label=right:$z_1$] (z1_right) at (4.2, 1.5) {};
         \node[circle, fill, inner sep=1.2pt, label=right:$z_2$] (z2_right) at (4.8, 0.5) {};
         \node[circle, fill, inner sep=1.2pt, label=right:$z_3$] (z3_right) at (4.0, -0.5) {};
         \node[circle, fill, inner sep=1.2pt, label=right:$z_4$] (z4_right) at (4.5, -1.5) {};

         % Nối mũi tên minh họa sự di chuyển của phần tử (giống style của bạn)
         \draw[->, thick, colorEmphasisCyan, shorten >= 5pt, shorten <= 5pt, dashed] (z1_left) to[bend left=15] (z1_right);
         \draw[->, thick, colorEmphasisCyan, shorten >= 5pt, shorten <= 5pt, dashed] (z3_left) to[bend right=15] (z3_right);

      \end{tikzpicture}
   \end{center}
   \caption{Minh họa tiên đề hợp (các phần tử của $A$ và $B$ trở thành phần tử trực tiếp của $\bigcup x$)}
\end{figure}

Dễ dàng nhận ra rằng tập này là duy nhất. Giả sử có hai tập $u_1$ và $u_2$ đều chứa tất cả và chỉ chứa các phần tử của các tập hợp trong $x$, khi đó ta có:
\begin{equation*}
   \begin{cases}
      \forall y \left( y \in u_1 \iff \exists z (z \in x \land y \in z) \right) \\
      \forall y \left( y \in u_2 \iff \exists z (z \in x \land y \in z) \right)
   \end{cases}.
\end{equation*}
Từ đó suy ra $\forall y, (y \in u_1 \iff y \in u_2)$, tức là $u_1 = u_2$.

Tập hợp này được gọi là \defText{hợp} của $x$ và kí hiệu là $$\defMath{\bigcup x}\text{ hoặc }\defMath{\bigcup_{A \in x} A}.$$

Đặt $x = \{A; B\}$. Khi này, mọi phần tử của A và B đều là phần tử của $\bigcup x$ và ngược lại. Chúng ta có thể hiểu $\bigcup x$ như là hợp của hai tập hợp A và B, tức định nghĩa tương đương \begin{equation*}
   \defMath{A \cup B = \bigcup \{A; B\}}.
\end{equation*}

\exercise Cho $A$, $B$, $C$ là các tập hợp. Chứng minh rằng
\begin{multicols}{2}
   \begin{enumerate}
      \item $\bigcup \{A\} = A$;
      \item $\bigcup \emptyset = \emptyset$;
      \item $A \cup B \supseteq A$;
      \item $\bigcup (A; B) = \{A; B\}$;
      \item $A = B \implies A \cup C = B \cup C$;
      \item $(A \cup B) \cup C = A \cup (B \cup C) = \bigcup \{A; B; C\}$.
   \end{enumerate}
\end{multicols}

\solution

\setcounter{subexercise}{1}
\arabic{subexercise}. Đặt $u = \bigcup \{A\}$. Gọi $x$ là một phần tử bất kì thuộc $A$. Hiển nhiên $A \in \{A\}$. Kết hợp với $x \in A$ để có
$$
   \exists z, (z \in \{A\} \land x \in z).
$$
Điều trên là đúng với mọi $x$. Từ đó, suy ra $\forall x, (x \in A \implies x \in u)$, dẫn đến $A \subseteq u$.

Ngược lại, gọi $y$ là một phần tử bất kì thuộc $u$. Theo định nghĩa của tiên đề hợp, chúng ta có
$$
   \exists z, (z \in \{A\} \land y \in z).
$$
Tuy nhiên, $\{A\}$ là một tập hợp chỉ có một phần tử duy nhất là $A$. Do đó, $z = A$ và $y \in A$. Tóm lại, với mọi $y$, nếu $y \in u$ thì $y \in A$, cho chúng ta $u \subseteq A$.

Từ hai điều trên, chúng ta có điều phải chứng minh.

\stepcounter{subexercise}
\arabic{subexercise}. Chúng ta chứng minh bằng phương pháp phản chứng. Giả sử $\bigcup \emptyset \neq \emptyset$, mà $x \notin \emptyset$ với mọi $x$, khi đó phải tồn tại ít nhất một phần tử $x \in \bigcup \emptyset$. Theo định nghĩa của hợp, sẽ tồn tại một $y \in \emptyset$ sao cho $x \in y$. Tuy nhiên điều này là vô lý, vì tập hợp rỗng $\emptyset$ không có chứa phần tử nào. Sự mâu thuẫn này chứng tỏ giả thiết phản chứng là sai. Vậy $\bigcup \emptyset = \emptyset$. Chúng ta có điều phải chứng minh.

\stepcounter{subexercise}
\arabic{subexercise}. Gọi $x \in A$. Nhắc lại rằng, $A \cup B = \bigcup \{A; B\}$. Khi đó, có $A \in \{A; B\}$ và giả sử ban đầu của $x$ để có $x \in A \cup B$ với mọi $x$. Vậy $A \cup B \supseteq A$. Điều phải chứng minh.

\stepcounter{subexercise}
\arabic{subexercise}. Trước tiên, nhắc lại định nghĩa đôi có thứ tự: $(A; B) = \big\{ \{A\}; \{A; B\} \big\}$. Lấy tùy ý $x \in \bigcup (A; B)$. Theo định nghĩa của hợp, tồn tại tập hợp $y \in \big\{ \{A\}; \{A; B\} \big\}$ sao cho $x \in y$. Do đó, $y = \{A\}$ hoặc $y = \{A; B\}$.
\begin{itemize}
   \item \textcolor{colorEmphasis}{Trường hợp $y = \{A\}$}: Chúng ta có $x \in \{A\}$ kéo theo $x = A$. Dẫn đến $x \in \{A; B\}$.
   \item \textcolor{colorEmphasisCyan}{Trường hợp $y = \{A; B\}$}: Chúng ta có ngay $x \in \{A; B\}$.
\end{itemize}
Cả hai trường hợp đều cho $\bigcup (A; B) \subseteq \{A; B\}$.

Ngược lại, lấy tùy ý $x \in \{A; B\}$. Có $\{A; B\} \in \big\{\{A\}; \{A; B\}\big\}$, nên $x \in \bigcup (A; B)$. Do vậy, $\{A; B\} \subseteq \bigcup (A; B)$.

Từ hai điều trên, chúng ta có điều phải chứng minh.

\stepcounter{subexercise}
\arabic{subexercise}. Giả sử $A = B$. Lấy tùy ý $x \in A \cup C$, theo định nghĩa, tồn tại tập hợp $y \in \{A; C\}$ sao cho $x \in y$. Khi này, $\left[\begin{array}{l}
      y = A \\
      y = C
   \end{array}\right.$.
\begin{itemize}
   \item \textcolor{colorEmphasis}{Trường hợp $y = A$}: Chúng ta có $x \in A$. Vì $A = B$, suy ra $x \in B$. Lại có $B \in \{B; C\}$, nên kéo theo $x$ là một phần tử của $B \cup C$.
   \item \textcolor{colorEmphasisCyan}{Trường hợp $y = C$}: Chúng ta có $x \in C$. Tương tự, có $x \in B \cup C$.
\end{itemize}
Cả hai trường hợp đều cho kết quả $x \in B \cup C$. Do đó, $A \cup C \subseteq B \cup C$.

Bằng lập luận tương tự, do vai trò của $A$ và $B$ là như nhau, chúng ta cũng xuất phát được với $x \in B \cup C$ và nhận về $B \cup C \subseteq A \cup C$. Từ hai kết quả trên, suy ra $A \cup C = B \cup C$. Điều phải chứng minh.
